\section{审核历史是怎样产生的}
\label{sec:introduction}
\subsection{从教材题目到经过审核的 Lean 文件}
\label{subsec:workflow}
这个项目把概率教材中的定义、定理、例题和习题写成 Lean 文件。Lean 是一种编程语言，也是一套检查形式化证明的系统。对于某个教材任务，编写智能体根据题目文字及相关数学材料写出候选文件，先检查它能否构建，再由审核智能体对照原题和所提供的审核要求，检查这个送审版本。审核者留下判断，并在适用时记录供编写者修改的问题。编写者随后可以提交另一个版本，再次接受检查和审核。

审核所问的不只是编译器是否接受文件，还包括：形式定义或陈述是否表达了题目要求的数学，证明是否处理了应有的论证，依赖和假设是否符合适用规则。因此，一次失败可能促使编写者修改陈述、证明或辅助定义；后来的审核也可能重新检查已经接受的文件。这些活动留下了不同材料：题目文字、保存的候选版本、构建结果、审核输入与要求、判断标签，以及问题说明。

这些记录首先服务于实际工作，随后才成为研究材料。编写者要知道改哪里，审核者要知道正在看哪道题的哪个版本，后来的维护者则要理解某个结果为什么曾被接受，又为什么重新审核。保留旧证据，才能把一次判断与当时可用的材料联系起来。保存被拒绝的候选或较早的审核，并不表示它仍是当前成果；这些记录也不是每一次编辑、每一次模型调用的完整日志。

以上是保留的流程文档和历史材料所支持的共同工作模式。档案还包括对正式成果的复查、外部成果审核，以及同一次审核的不同文件表示；要求、工具和记录方式都曾变化。因此，不能假定每条历史记录都由今天的形式化系统按照同一套实现产生。附录~\ref{sec:sources}~列出了这段背景说明所依据的材料。

\begin{figure}[htbp]
\centering\small
\fbox{\parbox{0.90\linewidth}{\textbf{教材题目与上下文。}确定要形式化的数学陈述、定义或问题。}}
\par\smallskip $\downarrow$ \par\smallskip
\fbox{\parbox{0.90\linewidth}{\textbf{候选与构建。}编写者提交保存的 Lean 版本，构建检查它在指定环境中是否被技术性地接受。}}
\par\smallskip $\downarrow$ \par\smallskip
\fbox{\parbox{0.90\linewidth}{\textbf{审核与下一步。}审核者核对送审数学，记录判断与问题。需要修改时回到候选步骤；接受之后还可能发生落地或后续复查。}}
\caption{反复出现的工作模式，证据随各项活动保留下来。图中并不假定每条历史记录都经历了每一个方框。}
\label{fig:work-pattern}
\end{figure}

\subsection{为什么从失败到通过仍需要解释}
一次提交涉及几个不同的问题：软件能否处理它？证明是否遵守指定依赖规则？形式陈述是否表达了原题的数学内容？这些问题的答案可能不同。例如，“假设结论成立，则结论成立”可以具有形式上有效的证明，却仍未解决原题。Lean 文档也区分了证明的有效性检查与陈述含义的核查\cite{leanvalidation}。

现在假设较早的审核拒绝了候选，较后的审核接受了它。编写者可能确实修正了数学问题，但题目要求或审核者掌握的证据也可能变了，先前的问题还可能没有再次被注意到。两个判断标签本身无法区分这些解释。历史值得研究，正是因为它不只保留最后一次通过：它还提供了一部分有关审核对象、检查和中间结果的证据。

\subsection{第二个流程：把工作记录整理成分析数据}
\label{subsec:analysis-workflow}
本研究从编写和审核活动结束后开始。首先整理保留证据，把它们与任务和候选标识关联起来，登记可以取得哪些审核材料与上下文。接着，沿用档案规定的同任务相邻配对，判断哪些配对能确定时间方向，再分类每条连接表示什么活动。候选改变、正式成果复查和同次审核的第二种文件表示，分别保留自己的含义。完成分类以后，才把符合条件的修改连接拼成连续历史。

由此产生两个历史问题。把一次选定修改的前后端点配对，可以询问记录判断怎样变化；从一次明确失败开始沿连续历史观察，则可以询问哪一步首次记录通过，或在哪里结束了观察而未见通过。它们来自同一档案，却使用不同的数据构造。第~\ref{sec:materials}~节会先解释单位、筛选规则和计数，再使用统计结果。

研究还用明确规定的检查考察部分保存候选，并使用人为构造的比较，以及一个可以穷尽全部输入的小任务。这些材料回答另一类问题：评价器究竟能提供候选的哪些信息？据此，报告研究三个相互关联的问题：
\begin{enumerate}[leftmargin=*]
\item \textbf{测量：}一项指定检查确认了什么，又留下了什么没有解决？
\item \textbf{比较：}在两条明确标准下分别比较两个候选，哪些量发生了变化？
\item \textbf{记录过程：}判断怎样在选定修改前后变化，又怎样分布在保留的失败起点历史中？
\end{enumerate}
即使缺少独立的数学正确性判断，我们仍能统计记录中的决定；此时，统计解释必须保持在这些决定的范围内。

\subsection{阅读顺序与定义查找}
章节沿用便于查阅的顺序：对象及其选取方式、评价量、数学论证、可执行证据、配对统计，以及观察到的过程。定义规定各项分析所用的对象和量，命题说明明确假设能够推出什么。说明性例子用于解释计算，不是额外实验。

第~\ref{sec:materials}~节定义任务、候选、记录、连接和连续片段。第~\ref{sec:framework}~节介绍分数，以及把两个候选放在两条标准下分别检查的四格比较。第~\ref{sec:mathematics}~节给出测量、停止和识别方面的论证，第~\ref{sec:audit}~节说明保留的可执行检查。第~\ref{sec:statistics}~节和第~\ref{sec:process}~节展开配对与过程统计。附录~\ref{sec:definition-index}~提供带页码的定义索引，附录~\ref{sec:sources}~将分析连接到随附文件。

报告假定读者掌握基础微积分、概率、线性代数、统计推断及其常用记号，在需要时介绍项目专用术语和统计构造。每个平均数都说明什么对象获得权重，每个比例都交代分母数的是什么。
